\begin{ouverture}[calcul]{6}{Les fractions, du tout au nombre}{Partie 2 · Fractions et nombres décimaux}
{Parts · Demi-droite · Entiers cachés · Encadrements}
\objectif{Lire, représenter et écrire des fractions, au-delà de 1.}
\objectif{Passer d’une fraction à l’écriture entier plus fraction.}
\objectif{Placer et repérer des fractions sur une demi-droite.}
\prerequis{Fractions du CM1 : sens, demi-droite, encadrements.}
\end{ouverture}
\pagestyle{otMain}\essentiel\fontsize{11}{15}\selectfont
\phantomsection\label{cours-6-1}
\essSection{1}{Des parts de toutes tailles}
\begin{definition}La fraction $\frac{7}{12}$ désigne 7 parts d’une unité partagée en 12 parts égales. Le dénominateur nomme le partage, le numérateur compte les parts. Au CM2, les partages vont jusqu’aux soixantièmes.\end{definition}
\begin{exemple}Une heure partagée en 60 minutes : chaque minute est $\frac1{60}$ d’heure, et 45 minutes représentent $\frac{45}{60}$ d’heure, soit trois quarts d’heure.\end{exemple}
\begin{propriete}Une fraction dépasse 1 quand son numérateur dépasse son dénominateur : $\frac{19}6=3+\frac16$, car $18$ sixièmes font 3 unités. L’écriture entier plus fraction situe immédiatement le nombre.\end{propriete}
\begin{savoirfaire}[Encadrer entre deux entiers]Encadrer $\frac{47}{10}$ : $\frac{40}{10}=4$ et $\frac{50}{10}=5$, donc $4<\frac{47}{10}<5$. On cherche les multiples du dénominateur voisins du numérateur.\end{savoirfaire}
\clearpage
\phantomsection\label{cours-6-2}
\essSection{2}{Sur la demi-droite graduée}
\begin{propriete}Les fractions sont des nombres : sur une demi-droite dont l’unité est partagée en 6, le point $\frac{19}6$ se place un pas après 3. Réciproquement, un point se repère par une fraction ou par un entier plus une fraction.\end{propriete}
\label{fig-cours-6-2-cours-droite-sixiemes}\figureGeo{cours-droite-sixiemes}

\begin{exemple}Sur une même droite graduée en douzièmes, $\frac6{12}$ tombe sur le même point que $\frac12$ : deux écritures peuvent repérer le même nombre.\end{exemple}
\begin{savoirfaire}[Placer vite et bien]Pour placer $2+\frac56$ : aller à l’entier 2, puis avancer de 5 pas de un sixième. Pour lire un point : compter les unités entières, puis les pas restants.\end{savoirfaire}

\clearpage\rubrique{Découverte}
\begin{blocLivre}{ex-6-1}\begin{exercice}[1]{Lire des parts}
Quelle fraction de l’heure représentent 7 minutes ? Et 45 minutes ?
\end{exercice}\end{blocLivre}
\begin{blocLivre}{ex-6-2}\begin{exercice}[1]{Représenter}
La bande est partagée en 12 parts égales. Colorier $\frac5{12}$ de la bande.
\label{fig-ex-6-2}\figureGeo{frac-2}
\end{exercice}\end{blocLivre}
\begin{blocLivre}{ex-6-3}\begin{exercice}[1]{Entier caché}
Écrire $\frac{19}6$ et $\frac{48}8$ plus simplement, avec un entier et, si besoin, une fraction.
\end{exercice}\end{blocLivre}
\begin{blocLivre}{ex-6-4}\begin{exercice}[1]{Recomposer}
Écrire $4+\frac3{10}$ puis $2+\frac56$ sous la forme d’une seule fraction.
\end{exercice}\end{blocLivre}
\begin{blocLivre}{ex-6-5}\begin{exercice}[1]{Encadrer}
Encadrer entre deux entiers consécutifs : $\frac{47}{10}$, puis $\frac{23}4$.
\end{exercice}\end{blocLivre}

\rubrique{Entraînement}
\begin{blocLivre}{ex-6-6}\begin{exercice}[2]{Lire la demi-droite}
L’unité est partagée en sixièmes. Lire les nombres repérés par A et B, de deux façons chacun.
\label{fig-ex-6-6}\figureGeo{fd-6}
\end{exercice}\end{blocLivre}
\begin{blocLivre}{ex-6-7}\begin{exercice}[2]{Placer des fractions}
Sur la demi-droite graduée en cinquièmes, placer $C\left(\frac85\right)$ et $D\left(2+\frac35\right)$.
\label{fig-ex-6-7}\figureGeo{fd-7}
\end{exercice}\end{blocLivre}
\begin{blocLivre}{ex-6-8}\begin{exercice}[2]{Le même point}
Sur une droite graduée en douzièmes, les fractions $\frac6{12}$ et $\frac12$ repèrent-elles le même point ? Justifier avec les pas.
\label{fig-ex-6-8}\figureGeo{fd-8}
\end{exercice}\end{blocLivre}
\begin{blocLivre}{ex-6-9}\begin{exercice}[2]{Un entier tout rond}
Compléter : $\frac{\square}{15}=2$ ; $\frac{36}{\square}=4$ ; $\frac{\square}{60}=1$.
\end{exercice}\end{blocLivre}
\begin{blocLivre}{ex-6-10}\begin{exercice}[2]{Les tours de piste}
Un coureur a parcouru $\frac{11}4$ de tour de piste. Combien de tours complets a-t-il bouclés ? Quelle fraction de tour lui reste-t-elle pour finir le troisième tour ?
\end{exercice}\end{blocLivre}

\rubrique{Problèmes}
\begin{blocLivre}{ex-6-11}\begin{exercice}[3]{Les gobelets}
Avec un pichet de $\frac{13}4$ de litre, combien de gobelets d’un quart de litre peut-on remplir ? Que reste-t-il après avoir rempli 3 gobelets ?
\end{exercice}\end{blocLivre}
\begin{blocLivre}{ex-6-12}\begin{exercice}[3]{Entre deux entiers}
Trouver trois fractions de dénominateur 6 strictement comprises entre 2 et 3.
\end{exercice}\end{blocLivre}
\begin{blocLivre}{ex-6-13}\begin{exercice}[3]{Le pot de peinture}
Un peintre a utilisé $\frac{35}{10}$ de pot de peinture pour un chantier. Écrire cette quantité avec un entier, puis dire combien de pots entamés il a ouverts au minimum.
\end{exercice}\end{blocLivre}
\begin{blocLivre}{ex-6-14}\begin{exercice}[3]{Comparer à l’unité et à 2}
Classer ces fractions en trois groupes - plus petites que 1, entre 1 et 2, plus grandes que 2 : $\frac79$ ; $\frac{13}9$ ; $\frac{25}9$ ; $\frac99$.
\end{exercice}\end{blocLivre}
\begin{blocLivre}{ex-6-15}\begin{exercice}[3]{Vrai ou faux ?}
Issa affirme : « $\frac{60}{59}$ et $\frac{59}{60}$ sont égales, elles utilisent les mêmes nombres. » Corriger et situer chaque fraction par rapport à 1.
\end{exercice}\end{blocLivre}

\clearpage
